% --- NUMBER SETS --------------------------------------------------------------

% inserts the symbol for the naturals
\newcommand{\N}{\mathbb{N}}

% inserts the symbol for the naturals with 0 included
\newcommand{\Nzero}{\mathbb{N}_0}

% inserts the symbol for the integers
\newcommand{\Z}{\mathbb{Z}}

% inserts the symbol for the rationals
\newcommand{\Q}{\mathbb{Q}}

% inserts the symbol for the real set
\newcommand{\R}{\mathbb{R}}

% inserts the symbol for the complex set
\newcommand{\C}{\mathbb{C}}

% inserts the symbol for the extended complex set
\newcommand{\ExtC}{\mathbb{C}_\infty}


% --- NUMBER SETS --------------------------------------------------------------

% inserts the symbol for the identity function
\newcommand{\Id}{\mathrm{I}}

% inserts the non italics Euler's number (roman e)
\newcommand{\E}{\mathrm{e}}

% inserts the non italics imaginary number (roman i)
\newcommand{\I}{\mathrm{i}}


% --- ONE VARIABLE DIFFERENTIATION ---------------------------------------------

% inserts the roman d for differentiation
\newcommand{\D}{\mathrm{d}}

% inserts the fraction derivative symbol, it needs two arguments:
% 1. the function to differentiante (symbol in the numerator)
% 2. the variable respect to which to differentiate (symbol in the denominator)
\newcommand{\dbyd}[2]{\frac{\D#1}{\D#2}}

% inserts the fraction n-th derivative symbol, it needs two arguments:
% 1. the function to differentiante (symbol in the numerator)
% 2. the variable respect to which to differentiate (symbol in the denominator)
% and can recieve an optional argument to indicate the order of the derivative
\newcommand{\dnbyd}[3][]{\frac{\D^{#1}#2}{\D #3^{#1}}}


% --- TOPOLOGY -----------------------------------------------------------------

% applies the interior operator
\newcommand{\interior}[1]{{#1}^{\mathrm{o}}}

% applies the closure operator
\newcommand{\closure}[1]{\overline{#1}}

% applies the boundary operator
\newcommand{\boundary}[1]{\partial {#1}}

% denotes the derived set
\newcommand{\derived}[1]{{#1}'}

% --- COMPLEX DYNAMICS ---------------------------------------------------------

% inserts the symbol for the orbit of a point
\newcommand{\orbit}{\mathcal{O}}